%% Canonical chain roster per LP-134 (Lux Chain Topology): the NINE chains
%% P, C, X, Q, Z, A, B, M, F. T-Chain is REMOVED entirely — its concerns
%% decompose with zero remainder onto M-Chain (mpcvm, all threshold signing),
%% F-Chain (fhevm, all FHE/confidential compute), B-Chain (bridgevm, teleport /
%% cross-chain messaging), Q-Chain (consensus-threshold), and Z-Chain (ZK).
%% thresholdvm is a shared LIBRARY substrate consumed by M and F, not a chain.
%% Teleport IS bridgevm; there is no teleportvm. See LP-7050 for the full map.
\documentclass[11pt,a4paper]{article}
\usepackage{shared/luxcover}
\usepackage[utf8]{inputenc}
\usepackage[T1]{fontenc}
\usepackage{amsmath,amssymb,amsthm}
\usepackage{graphicx}
\usepackage{hyperref}
\usepackage{booktabs}
\usepackage{multirow}
\usepackage[margin=1in]{geometry}

\hypersetup{colorlinks=true,linkcolor=black,citecolor=black,urlcolor=black}

\newtheorem{theorem}{Theorem}
\newtheorem{definition}{Definition}

\title{\textbf{Lux Network: A Multi-Chain Architecture\\for Orthogonal, Composable Blockchain Infrastructure}}

\author{Lux Research\thanks{research@lux.network}\\
\textit{Lux Industries}}

\date{February 2026}

\begin{document}
\luxcoverpage

\begin{abstract}
We present the complete multi-chain architecture of Lux Network. Its canonical roster (LP-134) is the nine chains \textbf{P, C, X, Q, Z, A, B, M, F}: the primary network (P, X, C) provides the consensus foundation, while Q, Z, A, B, M, and F each handle exactly one concern with no functional overlap. M-Chain (\texttt{mpcvm}) owns all threshold signing for bridge custody; F-Chain (\texttt{fhevm}) owns all FHE / confidential compute; both consume the shared \texttt{thresholdvm} library substrate. Teleport is the Bridge chain (\texttt{bridgevm}); there is no separate T-Chain or teleportvm. This paper additionally documents the operational D (DEX), I (Identity), K (Keys), and G (read-only query) chains that round out a running deployment. Cross-chain composition occurs via Warp messages at the consensus level~\cite{lux-bridge}. The D-Chain DEX~\cite{lux-dex} serves as a universal price oracle enabling gas payment in any listed asset across all chains. Quasar consensus~\cite{lux-consensus} coordinates C+Q+M+Z for post-quantum triple-signature receipts. This paper defines each chain's state model, transaction types, fee structure, and inter-chain composition protocol.
\end{abstract}

\section{Design Principles}

\begin{enumerate}
\item \textbf{One and one way only.} Each concern has exactly one chain. No redundancy.
\item \textbf{Separation of concerns.} Each chain owns its domain completely.
\item \textbf{Orthogonal.} Chains do not share mutable state. Changes on one cannot corrupt another.
\item \textbf{Composable.} Any chain can invoke any other chain's capabilities via Warp messages.
\item \textbf{Complete.} Every operation in the ecosystem has a canonical home chain.
\end{enumerate}

\section{Chain Registry}

\subsection{Primary Network}

\begin{table}[h]
\centering
\begin{tabular}{clllp{6cm}}
\toprule
\textbf{ID} & \textbf{Chain} & \textbf{VM} & \textbf{Gas} & \textbf{State} \\
\midrule
P & Platform & PlatformVM & LUX & Validator set, staking balances, L1 chain configs, delegation \\
X & Exchange & XVM & LUX & UTXO set, atomic swap state, NFT ownership \\
C & Contract & EVM & LUX\textsuperscript{$\dagger$} & Accounts, contract storage, precompile state \\
\bottomrule
\end{tabular}
\caption{Primary network chains. $\dagger$: C-Chain accepts any DEX-listed asset as gas payment.}
\end{table}

\subsection{Specialized Chains}

\begin{table}[h]
\centering
\begin{tabular}{clllp{5.5cm}}
\toprule
\textbf{ID} & \textbf{Chain} & \textbf{Concern} & \textbf{Gas} & \textbf{State} \\
\midrule
A & Attestation & TEE proofs & LUX & Attestation records, compute proofs, NVTrust certificates \\
B & Bridge & Cross-chain & LUX & Pending transfers, MPC signer sets, liquidity pools, route state \\
D & DEX & Trading & LUX\textsuperscript{$\dagger$} & Orderbook, pool state, positions, margin accounts \\
I & Identity & DIDs/VCs & LUX & DID documents, credential registries, revocation lists \\
K & Keys & Key mgmt & LUX & Public key records, rotation history, policy rules \\
M & MPC & Threshold signing & LUX & DKG session state, key shares, threshold group configs (bridge custody of external wallets; consumes \texttt{thresholdvm} library) \\
F & FHE & Confidential compute & LUX & TFHE bootstrap keys, encrypted EVM state, threshold-decrypt requests (consumes \texttt{thresholdvm} library) \\
Q & Quantum & Post-quantum & LUX & Pulsar key state, NTT tables, PQ sig aggregation \\
Z & ZK & Proofs & LUX & Proof registry, receipt Merkle tree, circuit registry \\
\bottomrule
\end{tabular}
\caption{Specialized consensus chains. $\dagger$: D-Chain fees paid in the traded pair.}
\end{table}

\subsection{Read-Only Layer}

\begin{table}[h]
\centering
\begin{tabular}{clllp{6cm}}
\toprule
\textbf{ID} & \textbf{Chain} & \textbf{Concern} & \textbf{Gas} & \textbf{State} \\
\midrule
G & Graph & Query & None & Read-only index of all chains. No consensus participation. Can run anywhere. \\
\bottomrule
\end{tabular}
\caption{Read-only query layer.}
\end{table}

\section{Universal Gas Payment via D-Chain}

Every Lux chain denominates fees in LUX but accepts any asset listed on the D-Chain DEX:

\begin{definition}[Universal Gas Payment]
For a transaction on chain $\mathcal{C}$ with gas cost $g$ LUX and a user holding asset $A$:
\begin{enumerate}
\item Fee manager queries D-Chain for the $A$/LUX spot price $p$.
\item User's $A$ balance is debited $g \cdot p$ units of $A$.
\item Validator receives $g$ LUX equivalent.
\item The operation is atomic---no separate swap transaction visible to the user.
\end{enumerate}
\end{definition}

Accepted gas assets include but are not limited to: LUX, LUSD, LETH, LBTC, ZOO, PARS, and any ERC-20 with a D-Chain liquidity pool.

\section{Sovereign L1 Integration}

Sovereign L1 chains (Zoo, Pars, Hanzo, custom L1s) participate in the architecture as follows:

\begin{itemize}
\item Each L1 has its own native coin (ZOO, PARS, AI, etc.)
\item L1 validators may be a subset of Lux primary validators or independent
\item Cross-L1 communication uses Warp messages via B-Chain
\item L1 coins are tradeable on D-Chain for gas flexibility
\item D-Chain oracle provides the price feed for cross-L1 gas payment
\end{itemize}

\section{Quasar Consensus}

The Quasar consensus protocol~\cite{lux-consensus} spans four chains for post-quantum block finality:

$$\text{Quasar} = \text{ZKP}\left(\text{BLS}_{C} \parallel \text{Pulsar}_{Q} \parallel \text{ML-DSA}_{K}\right)$$

\begin{enumerate}
\item \textbf{C-Chain}: Validators produce BLS12-381 aggregate signatures (fast path).
\item \textbf{Q-Chain}: Validators produce Pulsar LWE threshold signatures (PQ path).
\item \textbf{K-Chain}: ML-DSA keys sign the block hash (NIST PQ standard path).
\item \textbf{M-Chain}: Coordinates threshold signing across all three protocols.
\item \textbf{Z-Chain}~\cite{lux-zk}: Wraps all three verifications in a STARK receipt. The receipt is the proof of finality.
\end{enumerate}

\begin{theorem}[Quasar Finality]
A block is final when the Z-Chain receipt attests to valid signatures from all three paths with $\geq 67\%$ validator agreement on each. Under the assumption that at least one of BLS12-381, Corona-LWE, or ML-DSA~\cite{lux-pq} remains secure, finality cannot be forged.
\end{theorem}

\section{Fee Schedule}

\begin{table}[h]
\centering
\begin{tabular}{llrl}
\toprule
\textbf{Chain} & \textbf{Operation} & \textbf{Base Fee} & \textbf{Unit} \\
\midrule
P & Create validator & 0 & LUX \\
P & Add delegator & 0 & LUX \\
X & Transfer & 0.001 & LUX \\
C & EVM transaction & dynamic & LUX (EIP-1559) \\
A & Submit attestation & 0.01 & LUX \\
B & Initiate bridge transfer & 0.05 & LUX \\
D & Place order (maker) & 0 & traded pair \\
D & Fill order (taker) & 0.001\% & traded pair \\
D & Add liquidity & 0.01 & LUX \\
I & Register DID & 0.1 & LUX \\
I & Issue credential & 0.01 & LUX \\
K & Register public key & 0.05 & LUX \\
K & Rotate key & 0.02 & LUX \\
M & DKG session & 1.0 & LUX \\
M & Threshold sign & 0.1 & LUX \\
Q & Pulsar verify & 0.5 & LUX \\
Q & NTT batch & 0.2 & LUX \\
Z & STARK verify & 0.5 & LUX \\
Z & Groth16 verify & 0.2 & LUX \\
Z & Issue receipt & 0.1 & LUX \\
G & Query & 0 & free (read-only) \\
\bottomrule
\end{tabular}
\caption{Fee schedule by chain and operation.}
\end{table}

\section{ZapDB Storage}

Each chain stores its state in ZapDB, Lux's native embedded database:

\begin{itemize}
\item \textbf{P-Chain}: Validator records, delegation state, L1 chain membership (UTXO model)
\item \textbf{X-Chain}: UTXO set with asset metadata
\item \textbf{C-Chain}: MPT state trie (Ethereum-compatible)
\item \textbf{A-Chain}: Attestation Merkle tree (append-only)
\item \textbf{B-Chain}: Transfer queue + signer set + liquidity accounting
\item \textbf{D-Chain}: Orderbook B-tree + AMM pool state + position trie
\item \textbf{I-Chain}: DID document store + credential index + revocation bitmap
\item \textbf{K-Chain}: Key registry trie + policy rules + rotation log
\item \textbf{M-Chain}: Session state + share commitments + group configs
\item \textbf{Q-Chain}: NTT coefficient tables + Pulsar aggregation buffer
\item \textbf{Z-Chain}: Receipt Merkle tree + circuit registry + proof cache
\item \textbf{G-Chain}: Read-only aggregate index (no consensus writes)
\end{itemize}

\section{Conclusion}

Lux Network's multi-chain architecture (canonical roster: the nine chains P, C, X, Q, Z, A, B, M, F per LP-134) provides complete separation of concerns with orthogonal, composable chains. Each chain owns exactly one domain, stores its own state in ZapDB, defines its own transaction types, and sets its own fee schedule. The D-Chain DEX enables universal gas payment in any listed asset. Quasar consensus coordinates C+Q+M+K+Z for post-quantum finality. G-Chain provides a free, read-only query layer that requires no consensus participation. Warp messages compose chains at the protocol level without shared mutable state.

\begin{thebibliography}{10}
\bibitem{lux-consensus} Kelling, Z., ``Quasar Hybrid Consensus on Lux Network,'' Lux Industries, 2024.
\bibitem{lux-pq} Kelling, Z., ``Post-Quantum Cryptography Suite for Lux Network,'' Lux Industries, 2024.
\bibitem{lux-dex} Kelling, Z., ``Lux DEX: Native CLOB + AMM Precompiles,'' Lux Industries, 2024.
\bibitem{lux-bridge} Kelling, Z., ``Lux Bridge: Warp and Teleport Protocols,'' Lux Industries, 2024.
\bibitem{lux-zk} Worring, A., Kelling, Z., ``Z-Chain: Universal ZK Proof Platform,'' Zoo Labs Foundation, 2026.
\end{thebibliography}

\end{document}
